\section{Experiments}
\label{sec:experiments}
% 【中文】\section{实验}

{\hbadness=10000\sloppy \textbf{Experimental status.} We formulate 3 core hypotheses (H1--H3) and 7 sub-hypotheses (details in \S\ref{sec:hypotheses}). M0 is \emph{evaluated} on the full \Ncases{}-task set over 3 independent runs (\MzeroPass{}$\pm$\MzeroStd{} Pass@1; DeepSeek-V4-Pro, Cline v3.0.51, \texttt{thinking=low}, 2400s, 16 workers; r1=105/482, r2=96/486, r3=103/490; 3/3 clean, 0 FUSE/quota/burnout). M1--M8 and N1--N2 are to be evaluated. Negative results are reported as publishable boundary conditions.\par}
% 【中文】
% \textbf{实验状态。} 本文提出 3 条核心假设（H1--H3）与 7 条子假设（H1a--c、H2a--c、
% H3）。M0 基线已在完整 \Ncases{} 任务集上\emph{实测}（3 轮独立运行，
% \MzeroPass{}$\pm$\MzeroStd{} Pass@1；DeepSeek-V4-Pro、Cline v3.0.51、\texttt{thinking=low}、
% 2400s 超时、16 workers；r1=105/482、r2=96/486、r3=103/490；3/3 clean、0 FUSE/额度/熔断）。
% 所有其他臂（M1--M8、N1--N2）待评测。阴性结果将作为可发表的边界条件报告。

\textbf{M0 baseline \& backbone comparison.} Table~\ref{tab:m0-baseline} reports the M0 3-run snapshot. Table~\ref{tab:model-comparison} compares 4 backbones on the same \Ncases{}-task set (3 runs each, identical harness). The strongest backbone (DeepSeek-V4-Pro) reaches \MzeroPass{}$\pm$\MzeroStd{}, roughly 2.0$\times$ the weakest viable backbone (doubao-seed-2.0-pro-cloud, \DoubaoPro{}$\pm$\DoubaoProStd{}); a paired McNemar test on majority-vote per-task outcomes confirms the gap is highly significant ($b=69$, $c=2$, $p\approx 0.0000$, $\Delta{=}+10.67$pp).
% 【中文】
% \textbf{M0 基线与骨干对比。} 表~\ref{tab:m0-baseline} 报告 M0 的 3 轮快照。
% 表~\ref{tab:model-comparison} 在同一 \Ncases{} 任务集上对比 4 个骨干（各 3 轮、同一 harness）。
% 最强骨干（DeepSeek-V4-Pro）达到 \MzeroPass{}$\pm$\MzeroStd{}，约为最弱可用骨干
%（doubao-seed-2.0-pro-cloud，\DoubaoPro{}$\pm$\DoubaoProStd{}）的 2.0$\times$；
% 逐题 majority-vote 的 McNemar 配对检验确认差距高度显著
%（$b=69$、$c=2$、$p\approx 0.0000$、$\Delta{=}+10.67$pp）。

\begin{table}[ht]
\centering
\caption{M0 baseline (3-run snapshot on the \Ncases{}-task set, DeepSeek-V4-Pro, 2026-11-24). Pass\% = PASS/(PASS+FAIL). Full per-domain table in Appendix~\ref{sec:additional-results}.}
\label{tab:m0-baseline}
\scriptsize
\setlength{\tabcolsep}{4pt}
\begin{tabular}{@{}lcc@{}}
\toprule
\textbf{Metric} & \textbf{Count} & \textbf{Rate} \\
\midrule
PASS (3 runs total) & 304 & --- \\
FAIL (3 runs total) & 1154 & --- \\
ERROR (3 runs total) & 42 & --- \\
\midrule
Pass@1 (r1) & 105/482 & 21.78\% \\
Pass@1 (r2) & 96/486 & 19.75\% \\
Pass@1 (r3) & 103/490 & 21.02\% \\
Pass@1 (mean$\pm$std) & --- & \textbf{\MzeroPass{}$\pm$\MzeroStd{}} \\
\bottomrule
\end{tabular}
\end{table}

\begin{table}[ht]
\centering
\caption{Backbone comparison on the \Ncases{}-task set (3 runs each, same harness: Cline v3.0.51, \texttt{thinking=low}, 2400s, 16 workers). McNemar test on majority-vote per-task outcomes (DeepSeek-V4-Pro vs doubao-pro): $b=69$, $c=2$, $p\approx 0.0000$, $\Delta{=}+10.67$pp.}
\label{tab:model-comparison}
\scriptsize
\setlength{\tabcolsep}{4pt}
\begin{tabular}{@{}lccc@{}}
\toprule
\textbf{Backbone} & \textbf{Pass@1 (mean$\pm$std)} & \textbf{decided/run} & \textbf{clean} \\
\midrule
DeepSeek-V4-Pro & \textbf{\DSvPro{}$\pm$\DSvProStd{}} & 482/486/490 & 3/3 \\
doubao-seed-2.0-pro & \DoubaoPro{}$\pm$\DoubaoProStd{} & 483/483/484 & 3/3 \\
doubao-seed-2.0-lite & \DoubaoLite{}$\pm$\DoubaoLiteStd{} & 482/484/484 & 3/3 \\
doubao-seed-2.0-mini & \DoubaoMini{}$\pm$\DoubaoMiniStd{} & 483/483/484 & 3/3 \\
\bottomrule
\end{tabular}
\end{table}

\textbf{H1/H2/H3 criteria-check framework.} Each sub-hypothesis's pass/fail is determined by the pre-specified falsification threshold (\S\ref{sec:hypotheses}). Table~\ref{tab:criteria-check} gives the evaluation framework; results are \TBD{} (arms M1--M8, N1--N2 not yet evaluated).
% 【中文】
% \textbf{H1/H2/H3 判据检验框架。} 每条子假设的通过/失败由事前规定的证伪阈值（\S\ref{sec:hypotheses}）
% 决定。表~\ref{tab:criteria-check} 给出评测框架，结果为 \TBD{}（臂 M1--M8、N1--N2 尚未评测）。
% 【中文】表 3：判据检验（结果 \TBD{}，臂未评测）。

\begin{table}[ht]
\centering
\caption{Criteria check (results \TBD{}; arms M1--M8, N1--N2 pending evaluation).}
\label{tab:criteria-check}
\scriptsize
\begin{tabularx}{\columnwidth}{@{}llYcc@{}}
\toprule
\textbf{Hyp} & \textbf{Arm} & \textbf{Test} & \textbf{Est.} & \textbf{Verdict} \\
\midrule
H1a & M3 vs M0 & $\Delta$Pass@1 $\geq$3pp & \TBD & \TBD \\
H1a & M3 vs N1 & exclude prompt-lengthening & \TBD & \TBD \\
H1b & M3+S1 stratified & monotonic $\Delta$ & \TBD & \TBD \\
H1c & 2$\times$2 factorial & interaction $\geq$2pp & \TBD & \TBD \\
H2a & M5 vs M2 & surprise-gate $>$ uniform & \TBD & \TBD \\
H2b & M7 vs M2 & compress $>$ full & \TBD & \TBD \\
H2c & M6 vs M3 & state-aware $>$ agnostic & \TBD & \TBD \\
H3 & single vs spawn & single $\geq$ spawn & \TBD & \TBD \\
\bottomrule
\end{tabularx}
\end{table}

\textbf{Arm-table convergence \& power notes.} The main 10-arm table drops M1/M8 to appendix; M2 is the H2a/H2b control (H1a weak-control criteria withdrawn, \S\ref{sec:arms}). M6's fail-streak test is exploratory at the minimal matrix (H2c deferred to ideal matrix). The minimal-matrix M0 anchor is the 3-run M0 on the full \Ncases{}-task set (\MzeroPass{}$\pm$\MzeroStd{}), not a subset.
% 【中文】
% \textbf{臂表收敛与功效注。} 主 10 臂表将 M1（静态 SOP）与 M8（负效价）移至附录，M2 作为 H2a/H2b 对照
% 运行（非主臂；其 H1a 弱对照判据已撤回，§\ref{sec:arms}）。M6 的失败连续检验在最小矩阵下为探索性
%（H2c 判据延至理想矩阵）。最小矩阵的 M0 比较锚点为完整 \Ncases{} 任务集上的 3 轮 M0
%（\MzeroPass{}$\pm$\MzeroStd{}），而非子集。

\textbf{M0 measurement \& reused assets.} M0 is measured directly on the arm backbone DeepSeek-V4-Pro over the full \Ncases{}-task set: \MzeroPass{}$\pm$\MzeroStd{} Pass@1 (3 runs, r1=105/482, r2=96/486, r3=103/490, 2026-11-24). For H1b stratification, we re-use the trace-length distribution from 47 pilot runs (median 26.5 calls, IQR 13.8--38.0) to define short ($<$14)/mid (14--38)/long ($>$38) strata. M0 is evaluated; M1--M8 and N1--N2 are pending.
% 【中文】
% \textbf{M0 实测与复用资产。} M0 直接在臂骨干 DeepSeek-V4-Pro 上、于完整 \Ncases{} 任务集实测：
% \MzeroPass{}$\pm$\MzeroStd{} Pass@1（3 轮，r1=105/482、r2=96/486、r3=103/490，2026-11-24）。
% 对需要分层的臂（H1b），我们复用 47 条插桩试点 run 的轨迹长度分布
%（中位 26.5 调用，IQR 13.8--38.0）来定义短（$<$14）/中（14--38）/长（$>$38）分层。M0 已评测；
% M1--M8、N1--N2 待评测。
